\section{The cuspidal triangle endpoint}\label{sec:endpoint}

We complete the Cartan estimate at the endpoints of
Proposition~\ref{prop:runner-reduction}.  The ordinary cuspidal rank may
still be arbitrarily large there.  The number of split Harish--Chandra
steps above the cuspidal label is bounded, and this is the rank that will
enter the endomorphism algebras below.  We construct enough modular
cuspidal pairs to exhaust the simple modules of the endpoint blocks,
and then cover their projective indecomposable modules by bounded
projective inductions.

The use of intertwining operators and Hecke endomorphism algebras
follows the Harish--Chandra theory of Howlett--Lehrer
\cite{HowlettLehrer} and its modular developments
\cite{DipperFleischmann,DipperDu}.  We give the particular cuspidal,
extension and exhaustion arguments required here.

Throughout this section, $p$ is odd and different from the defining
characteristic $r$, and $k=\overline{\F}_p$.  We retain the regular
ambient groups, marked semisimple $p'$-parameter $s$, and actual Levi
product decompositions of the preceding sections.  Write
$G=\mathbf G^F$ and $\mathbf C=C_{\mathbf G^*}(s)$.  A packet means a
Frobenius orbit of classical factors of this connected centralizer on
which the labels vary.  We calculate on one representative factor
with the composite endomorphism around the orbit.  Index the packets
by $i$.  Their parameters satisfy
one of the following two conditions:
\begin{equation}\label{end:linear-primes}
\begin{array}{ll}
\text{orthogonal or symplectic:}&
 Q_i=q_i,\quad d_i=\ord_p(Q_i)\text{ is odd};\\
\text{unitary:}&
 Q_i=q^2,\quad d_i=\ord_p(Q_i),\quad \ord_p(-q)=2d_i.
\end{array}
\end{equation}
There is no assertion that $d_i$ is odd in the unitary case.
On each packet fix an ordinary cuspidal triangle: a separately packed
symbol in types $B,C,D$, or the staircase $t(t+1)/2$ in unitary type.
Let $m_i$ be the number of ordinary split Harish--Chandra steps above
it.  The collection $\mathcal T$ consists of all the ordinary rows at
$s$ with these triangle data and these $m_i$.  Frozen factors carry
their prescribed single cuspidal row and are left implicit.

For a type $D$ packet, zero difference between the two row lengths
means split type and an empty residual triangle.  We call this the
\emph{even-sign case}.  A nonempty type $D$ triangle has its unique
cuspidal label, also when its rank is one and its form is nonsplit.
At rank zero there is one label; in particular, the empty equal pair
is not doubled.  Equal nonempty pairs in split type $D$ retain their
two distinct labels.

\begin{proposition}[Triangle Cartan bound]\label{prop:triangle-cartan}
In the endpoint situation of Proposition~\ref{prop:runner-reduction},
the Cartan entries of the blocks whose basic rows belong to
$\mathcal T$ are bounded in terms of $p$ and the original defect-order
bound $M$.
More explicitly, the conclusion is uniform whenever the number of
packets, $\sum_i m_i$, the orders of the central $p$-subgroups retained
in the residual factors, and
$p^{v_p(Q_i^{d_i}-1)}$ for the packets with $m_i>0$ are bounded.
The bound is independent of the ranks of the triangles and of $q$.
\end{proposition}

The hypotheses listed in the second sentence are provided by
Proposition~\ref{prop:runner-reduction}.  That proposition and
Proposition~\ref{prop:core-blocks} also show that $\mathcal T$ is the
restriction of an integral basic set to a union of blocks of bounded
defect orders.  Both its cardinality and the number of ordinary rows
in this block union are bounded.  We will use these facts only at the
exhaustion and norm steps.  The modular Harish--Chandra calculation is
proved directly below.

\subsection{Residual and general linear cuspidal modules}

For each packet set
\begin{equation}\label{end:sizes}
 \mathcal B_i=\{1\}\cup\{d_i p^a:a\geq0\}.
\end{equation}
The sizes \(d_i p^a\) arise as Frobenius orbit degrees of regular
\(p\)-power eigenvalues over \(\F_{Q_i}\).  They will give ordinary
cuspidal lifts on the linear factors; size \(1\) also allows the
original \(p'\)-parameter with no added \(p\)-part.
This is a set: when $d_i=1$, size $1$ occurs only once.  Choose
nonnegative integers $r_{i,b}$ such that
\begin{equation}\label{end:size-sum}
 \sum_{b\in\mathcal B_i}b r_{i,b}=m_i.
\end{equation}
Partition the $m_i$ hyperbolic coordinates into these chunks.  In
$\mathbf C$ they define a split Levi with the residual triangle
factors and linear factors $\GL_b(Q_i)$.  Centralizing the same split
torus directions in $\mathbf G^*$ gives a rational split Levi
$\mathbf L^*$ with
\[
 C_{\mathbf L^*}(s)=\mathbf C_L.
\]
Write $L=\mathbf L^F$ for its dual Levi.  The actual product
description from Section~\ref{sec:families} identifies its varying
factors with a residual group $H_0$ and general linear groups
$\GL_{\delta_i b}(q)$ in the orthogonal and symplectic cases, where
$\delta_i\in\{1,2\}$ and $Q_i=q^{\delta_i}$, or $\GL_b(q^2)$ in
the unitary case.  On the former factor $s$ has one $p'$-eigenvalue
orbit $f_i$ of degree $\delta_i$, repeated $b$ times.  On the latter
factor it is the identity.

We record the split-center observation needed for these Levis.
The connected centralizer of the residual parameter has no split
central direction outside the ambient center.  A nonempty classical
triangle retains its full sign root system; the rank-one type $D$
exception is a nonsplit torus and has no such split direction either.
On a frozen linear factor with a regular torus parameter, the split
center of its centralizer is the split center of that ambient factor.
Consequently the split center of $\mathbf C_L$, modulo ambient
central directions, consists exactly of the chunk directions.  It
has the same centralizer as the split center of $\mathbf L^*$.
These assertions are statements about the rational cocharacter
spaces, so they are unaffected by the bounded central lattice
dressing.  All ensuing Levi conjugacies and Weyl positions are
rational.

\begin{lemma}[The residual module]\label{end:residual}
Let $\chi_0\in\Irr(H_0)$ have parameter $s$ and the indicated
cuspidal unipotent Jordan labels.  It is ordinary Harish--Chandra
cuspidal.  It is trivial on $Z_0=O_p(Z(H_0))$, and its character of
$H_0/Z_0$ has defect zero.  Its reduction $X_0$ is therefore simple
and Harish--Chandra cuspidal.  The block containing $X_0$ has bounded
Cartan data, and its only basic row at the residual parameter is
$\chi_0$.
\end{lemma}

\begin{proof}
Consider a proper rational split Levi of $H_0$ whose dual contains a
conjugate of the parameter.  Its intersection with the connected
centralizer is proper: otherwise its extra split central direction
would contradict the split-center observation.  Torus transitivity
and the ordinary Jordan torus pairing used in
Proposition~\ref{prop:single-cycle} identify the uniform projection
of the adjoint character there with the adjoint of the unipotent
cuspidal label.  The latter is zero.  Apply this separately to every
Levi parameter mapping to $s$.  In each such series the degree
vector belongs to the span of the torus tests, by the regular-character
formula.  The adjoint split Harish--Chandra character is an actual
character with nonnegative multiplicities; zero pairing with its
degree vector forces it to vanish.  This proves ordinary cuspidality.
No bound on the number or lengths of the torus tests is needed here.

The hook and cohook degree formula of
Proposition~\ref{prop:label-degrees} gives the defect assertion.
A separately packed triangle has no positive hooks.  Its cohook
factors have the form $Q_i^j+1$, whose $p$-parts are trivial because
$Q_i$ has odd order modulo the odd prime $p$.  The nonsplit rank-one
type $D$ torus contributes $Q_i+1$, also of order prime to $p$.
This accounts for the possible extra direction in the full residual
center: a rank-one orthogonal residual triangle is nonsplit, so
adjoining its torus does not enlarge the inherited central
$p$-subgroup.  The inherited-torus degree formula therefore leaves
the order of the full group $Z_0=O_p(Z(H_0))$.
Every hook length of a staircase is odd, whereas
$p\mid((-q)^j-1)$ requires $2d_i\mid j$.  Thus no noncentral
$p$-part remains in the character defect.  The order and degree
formulas for the regular ambient group leave precisely $|Z_0|$.
The central character attached to the $p'$-parameter is trivial on
$Z_0$, so passage to $H_0/Z_0$ is literal and gives defect zero.

A defect-zero ordinary character reduces to a simple projective
module.  Inflation from the central $p$-quotient gives the simple
module $X_0$, and the central transfer of
Lemma~\ref{lem:central-transfer} bounds its block's Cartan data in
terms of $|Z_0|$.  Exactness of split Harish--Chandra restriction,
which is averaging over groups of order prime to $p$, proves the
cuspidality of $X_0$.  Finally, every basic row in this block at the
specified $p'$-parameter is trivial on $Z_0$.  Downstairs it must lie
in the same one-row defect-zero block, and hence equals $\chi_0$.
\end{proof}

\begin{lemma}[Cuspidal modules on the chunks]\label{end:linear-cuspidals}
For every $b\in\mathcal B_i$ there is a simple cuspidal module
$Y_{i,b}$ on the corresponding actual general linear factor with
the following properties.
\begin{enumerate}
\item It is the reduction of an ordinary cuspidal character whose
semisimple parameter is the product of the prescribed $p'$-orbit
and a $p$-element regular of degree $b$ over $\F_{Q_i}$; for $b=1$
the added $p$-part may be trivial.
\item In the integral basic set of partition rows of the original
$p'$-series, its Brauer character is the signed power sum of degree
$b$.  In particular, it is independent of the choice of regular
$p$-part.
\item The ordinary partition row corresponding to the Steinberg
character on the unipotent side, called the regular extreme,
contains $Y_{i,b}$ with multiplicity one in its reduction.
This extreme is stable
under the packet transports and orientation reversals.
\end{enumerate}
\end{lemma}

\begin{proof}
Put $d=d_i$, $Q=Q_i$, and $c=v_p(Q^d-1)$.  Since $p$ is odd,
lifting the exponent gives
\[
 \ord_{p^{c+a}}(Q)=dp^a\qquad(a\geq0).
\]
For $b=dp^a>1$, an element $\mu$ of order $p^{c+a}$ therefore has
Frobenius orbit of length $b$ over $\F_Q$.  If the original orbit
has degree $\delta=2$, its product with $\mu$ has degree $2b$ over
$\F_q$.  Indeed its $p'$- and $p$-parts can be recovered separately,
so its orbit length is the least common multiple of their orbit
lengths.  The orbit length of $\mu$ over $\F_q$ is either $2b$, or
$b$ with $b$ odd, and either possibility gives least common multiple
$2b$ with $2$.  The case $\delta=1$ is immediate.  The resulting
torus parameter is in general position, so Green's ordinary
general linear theory gives its irreducible cuspidal character.
For $b=1$ use the ordinary cuspidal character of the original orbit.

Write $s_\lambda$ for the partition row with label $\lambda$, and
$p_b$ for the degree-$b$ power sum under Green's characteristic map.
The general linear torus rule and equality of torus characters on
$p$-regular elements give
\begin{equation}\label{end:power-sum}
 [\overline{\psi}_{i,b}]
   =\varepsilon_b p_b
   =\varepsilon_b\sum_{\lambda\vdash b}
       \chi^\lambda((b))[\overline{s_\lambda}],
 \qquad \varepsilon_b\in\{1,-1\}.
\end{equation}
Here $\varepsilon_b$ is the sign making the degree positive; in
degree one this is the single partition row.  The statements used
from Green's theory are the ordinary Coxeter-torus cuspidality,
this torus-character formula, and the identification of split
restriction with the symmetric-function coproduct~\cite{Green}.
For precise forms over the actual general linear factors, see
\cite[Property~(2.3a) and Lemma~2.3c]{BrundanDipperKleshchev}
for the power-sum identity on all $p$-regular classes, and
\cite[Equations~(2.2a), (2.3f) and Lemma~2.2d]{BrundanDipperKleshchev}
for induction and its adjoint restriction coproduct.

We give the irreducibility argument, since it avoids a modular
cuspidal classification.  Every composition factor $Y$ of
$\overline{\psi}_{i,b}$ has zero proper split restriction, by
exactness and ordinary cuspidality.  Express $[Y]$ integrally in
the partition basic set.  Its coproduct in every positive splitting
$b=b_1+b_2$ is zero, using the corresponding orbit-preserving Levi
and Green's restriction rule.  The primitive subspace of degree
$b$ in the rational symmetric-function Hopf algebra is $\Q p_b$:
indeed that algebra is the polynomial algebra on the primitive
generators $p_j$, and the reduced coproduct of a polynomial of
polynomial degree at least two is nonzero in characteristic zero.
The coefficient of an extreme Schur function in $p_b$ is $1$ or
$-1$, so $p_b$ is indivisible in the integral Schur lattice.  Thus
$[Y]=n_Y\varepsilon_b p_b$ for an integer $n_Y$.  Evaluation at the
identity shows $n_Y>0$.  Formula~\eqref{end:power-sum}, including
composition multiplicities, now gives $\sum_Y [\overline\psi_{i,b}:Y]
n_Y=1$.  There is exactly one factor, with multiplicity one.

For the extreme row, use the ordinary multiplicity-free
Gelfand--Graev formula for a general linear group
\cite[Theorem~2.5d(ii), (iv), (v)]{BrundanDipperKleshchev}.  Its
Gelfand--Graev character $\Gamma$ contains the regular Coxeter-torus
row $\psi_{i,b}$ once; among the partition rows at the original
parameter it contains only the regular extreme, namely the row
corresponding to Steinberg on the unipotent side, also once.
The Gelfand--Graev module is projective in characteristic $p$, since
it is induced from a linear character of a defining-characteristic
unipotent subgroup.  Let $P_Y$ be the projective cover of $Y_{i,b}$.
Brauer reciprocity and the irreducible reduction of $\psi_{i,b}$
show that $P_Y$ occurs once in this projective module.  Some basic
partition row has a nonzero decomposition number for $Y_{i,b}$,
because those rows span the Brauer lattice.  Nonnegativity then
forces that row to be the regular extreme and forces its
decomposition number to be one.  The extreme partition is carried
to the same extreme under all the relevant duality and field
transports.  This proves the last assertion.
\end{proof}

Form the simple cuspidal $L$-module
\begin{equation}\label{end:rho}
 \rho=X_0\otimes
       \bigotimes_{i,b}Y_{i,b}^{\otimes r_{i,b}},
\end{equation}
with the fixed cuspidal factors understood.  It has ordinary
cuspidal lifts $\psi^u$ with parameter $su$.  We may prescribe the
$p$-parts independently on the hyperbolic chunk coordinates of
$\mathbf C_L$.  In the orthogonal and symplectic ambient groups
the central and coordinate lattice gluing has only powers of $2$
as denominators and indices, hence is an isomorphism on $p$-power
torus data.  These prescriptions therefore give genuine rational
torus points.  Their projection to the residual factor can only
give a central $p$-power twist of $\chi_0$, which does not change
$X_0$.  The same construction in the unitary case is direct.

In normalizer notation below we always retain the algebraic Levi:
$N_G(\mathbf L,\rho)$ denotes the stabilizer of $\rho$ in
$N_{\mathbf G}(\mathbf L)^F$.  This convention is relevant in small
fields, where the abstract normalizer of the finite group $L$ need
not determine the algebraic Levi.

\begin{lemma}[Inertia]\label{end:inertia}
The stabilizer quotient $W_\rho=N_G(\mathbf L,\rho)/L$ is the relative
rational Weyl group of $\mathbf C_L$ in $\mathbf C$.  On a packet
outside the even-sign case its factor is
$\prod_b W(B_{r_{i,b}})$.  On an even-sign packet it is
\begin{equation}\label{end:total-parity}
 \left\{(w_b)_b\in\prod_b W(B_{r_{i,b}}):
          \sum_b \operatorname{flip}(w_b)=0\pmod2\right\},
\end{equation}
where $\operatorname{flip}$ is the parity of the number of changed
coordinate signs.  Different choices of the size multiplicities
$r_{i,b}$ give nonconjugate cuspidal pairs.
\end{lemma}

\begin{proof}
Relative normalizers on the group and dual group are identified by
their rational Weyl positions; Lang's theorem supplies rational
representatives.  An element preserving $\rho$ must preserve its
marked $p'$-series in $L$.  After adjustment by $\mathbf L^{*F}$,
its dual representative centralizes $s$ and hence belongs to
$\mathbf C^F$.  Conversely a rational normalizer of $\mathbf C_L$
in $\mathbf C$ preserves $\mathbf L^*$ by the split-center
description.  It preserves the residual cusp label and the power
sums in~\eqref{end:power-sum}; hence it fixes the Brauer character
of $\rho$.  The ordinary label transports here are exactly the
uniform and single-cycle transports of
Proposition~\ref{prop:single-cycle}.

In the hyperbolic coordinates, the resulting operations permute
chunks of equal size and reverse their orientations.  In an even
orthogonal factor a reversal of a chunk of size $b$ has coordinate
sign parity $b$.  All permitted sizes on that packet are odd,
since $d_i$ and $p$ are odd.  If the residual triangle is empty,
the condition is therefore the total parity in
\eqref{end:total-parity}.  It is a condition across all sizes, not
one condition for each size.  If the residual triangle is nonempty,
a single reversal can be compensated on it.  Its unique cusp label
is unchanged, including in the nonsplit rank-one case.  These
coordinate operations are Frobenius invariant and admit the
rational representatives just described.  No inner transporter in
the connected centralizer exchanges different eigenvalue packets.
Finally, the same transporter test preserves the multiset of chunk
sizes, proving nonconjugacy of the different displayed pairs.
\end{proof}

\subsection{Cuspidal induction and its heads}

The inducing modules and their inertia quotients are now explicit.
Before calculating the induction endomorphism algebras, we explain
what their simple types count.  Cuspidal Mackey theory identifies
them with the distinct simple modules in the heads of the induced
modules.  It also separates the heads attached to our nonconjugate
cuspidal pairs.

The Mackey calculation applies over $k$ or over a splitting field
of characteristic zero.
Write $R_L^G$ and ${}^*R_L^G$ for
split Harish--Chandra induction and restriction.  They are exact
biadjoint functors.  In the split Mackey formula, a double-parabolic
coset contributes restriction to an intersection Levi followed by
induction from it.  One obtains this formula by projecting the
parabolic intersection to its two Levi quotients and averaging
over the intersection unipotent groups.  Their orders are powers
of $r$ and hence invertible in either coefficient field.  With
cuspidal coefficients, proper intersection-Levi terms vanish.
Consequently
\begin{equation}\label{end:mackey}
 {}^*R_L^G R_L^G(\rho)
   \text{ is a direct sum of conjugate simple $L$-modules},
 \qquad
 \dim\End_G(R_L^G\rho)=|W_\rho|.
\end{equation}
Each normalizing Mackey position supplies a one-dimensional
intertwiner space.

\begin{lemma}[The head of cuspidal induction]\label{end:head}
Let $\rho$ be a cuspidal simple $L$-module and put
$V=R_L^G\rho$, $H=V/\rad V$.  Under the cuspidal Mackey
formula~\eqref{end:mackey}, the map
\[
 \End_G(V)\longrightarrow\End_G(H)
\]
is surjective with nilpotent kernel.  Its simple algebra
types therefore correspond to the distinct simple modules
in $H$.  Every simple quotient of $V$ is also a submodule
of $V$.  Heads arising from nonconjugate cuspidal pairs
have no common simple constituent.
\end{lemma}

\begin{proof}
Write $E={}^*R_L^G$, and let $\pi:V\to H$ be the quotient.
Exactness gives a surjection $E\pi:EV\to EH$.  The source
is semisimple by~\eqref{end:mackey}, so this surjection
splits.  Given $f\in\End_G(H)$, adjunction converts
$f\pi:V\to H$ into a map $\rho\to EH$, which lifts to
$\rho\to EV$.  Adjunction back supplies $g\in\End_G(V)$
with $\pi g=f\pi$.  This proves surjectivity.

Put $J=\rad(kG)$ and let $I$ be the kernel.  For $g\in I$,
$g(V)\subseteq JV$.  Since $g$ is $kG$-linear,
$g(J^aV)\subseteq J^{a+1}V$.  Thus a product of $n$
elements of $I$ maps $V$ into $J^nV$, which is zero for
large $n$.  The kernel is nilpotent.  Since $H$ is
semisimple, its endomorphism algebra is a product of full
matrix algebras, one per distinct simple constituent,
giving the first conclusion.

If $S$ is a simple quotient of $V$, then $ES$ is a
semisimple quotient of $EV$.  Adjunction provides
$\rho\hookrightarrow ES$, hence also $ES\twoheadrightarrow
\rho$ by semisimplicity.  The other adjunction gives a
nonzero map $S\to V$, which is an embedding.  If $S$
belongs to two such heads, composing a surjection from
one induced module with its embedding into the other
gives a nonzero homomorphism between the inductions.
The split Mackey formula and cuspidality force the two
pairs to be conjugate.  This proves disjointness.
\end{proof}

The pairs of Lemma~\ref{end:inertia} therefore give disjoint heads.
To prove that these heads exhaust the endpoint simples, we will
count their endomorphism-algebra types and show that the induced
modules are supported on the block union of $\mathcal T$.
The latter support assertion will be checked when the count is
complete.  We now determine the algebras whose simple types must
be counted.

\subsection{Chamber operators and removal of the scalar cocycle}

We calculate chamber operators and then extend the inducing module
to its inertia group to remove the scalar ambiguity in their
coefficient transports.  The operator calculation works over
either $k$ or a splitting field of characteristic zero.

Parabolics with Levi $\mathbf L$ are chambers on its split central
cocharacter space.  On the chunk coordinates their walls are among
$x_j=x_h$, $x_j=-x_h$, and $x_j=0$.  Add missing walls as dummy
walls so that the full signed-permutation arrangement is available.
Ambient central coordinates play no role.  If $U$ is the group of
rational points of a chamber's unipotent radical, put
\[
 e_U=|U|^{-1}\sum_{u\in U}u,
 \qquad
 \mathcal M_U=kG e_U\otimes_{kL}\rho.
\]
Choose a square root of $r$ and use its powers for all the following
normalizations.  Right multiplication by $e_V$ defines the
normalized change map
\begin{equation}\label{end:average}
 I_{U,V}:\mathcal M_U\longrightarrow\mathcal M_V,
 \qquad
 x e_U\otimes v\longmapsto
 \left(\frac{|V|}{|V\cap U|}\right)^{1/2}x e_U e_V\otimes v.
\end{equation}
It is well defined because $L$ normalizes both radicals.  The same
formula is used in characteristic zero.

\begin{lemma}[Averaging relations]\label{end:averaging}
Normalized changes along a minimal full gallery compose to the
direct map~\eqref{end:average}.  At an adjacent wall the reverse
composition is invertible.  If no inertial reflection fixes that
wall it is the identity.  Otherwise, after completing the change
by an involutively normalized coefficient transport for the
reflection $a$, the resulting endomorphism satisfies
\begin{equation}\label{end:quadratic}
 T_a^2=1+\alpha_a T_a.
\end{equation}
In particular $T_a^{-1}=T_a-\alpha_a$.
\end{lemma}

\begin{proof}
Let $V$ be an intermediate radical on a minimal gallery from $U$
to $U'$.  A root cannot cross its hyperplane twice on that gallery.
Thus the roots positive in $V$ are covered by those positive also
in $U$ or also in $U'$.  Set $A=V\cap U$ and $B=V\cap U'$.  Root
group order counting, including the overlap, gives
$|A||B|/|A\cap B|=|V|$.  Since $A,B\leq V$, this proves $AB=V$
as sets and $e_V=e_Ae_B$.  It follows that
\[
 e_Ue_Ve_{U'}=e_Ue_Ae_Be_{U'}=e_Ue_{U'}.
\]
This uses absorption of $e_A$ on the left and of $e_B$ on the
right, not commutation of arbitrary averaging idempotents.
The separating root sets are disjoint along a minimal gallery,
so the exponents in the square-root normalizations add.  This
proves the first assertion.

At a wall, first remove the common outer radical by transitivity.
Inside the Levi centralizing a generic point of the wall, the two
remaining radicals are opposite.  The coefficient of the identity
Mackey position in their reverse composition is $1$: an element of
the opposite radical contributes back to the identity parabolic
coset only if it is the identity, and the reciprocal radical order
is cancelled by the two normalizing factors.  If the wall has no
inertial reflection, the identity is the only surviving Mackey
position and gives the asserted inverse.

In the other case choose a representative $n_a$ and a coefficient
intertwiner whose square is the action of $n_a^2\in L$.  Such a
normalization is possible over a splitting algebraically closed
field by rescaling the intertwiner.  It is chosen for each reflection
separately and does not assume an extension to the full inertia group.
The square of the coefficient transport cancels right translation
by $n_a^{-2}$.  Right translation together
with that intertwiner completes the change to $T_a$.  Coefficient
transports commute covariantly with uncompleted averages.  Its
square is therefore the reverse composition.  Rank-one Mackey
has only the identity and reflection positions.  The reflection
position is nonzero: at its representative, the contributing
terms are precisely the intersection unipotents and carry the
same intertwiner.  The identity
$U\cap{}^{n_a}(LU)=U\cap{}^{n_a}U$, obtained from the root groups
of the common Levi, verifies this directly.  The normalized
identity coefficient is $1$, proving~\eqref{end:quadratic}.
\end{proof}

Let $W_{\rm r}$ be the reflection subgroup of $W_\rho$.  It is a
product of the indicated type $B$ groups, with
$\prod_b W(D_{r_{i,b}})$ on an even-sign packet.  Use
$W(D_1)=1$, $W(D_2)=C_2\times C_2$, and $W(D_0)=1$.
The subgroup $W_{\rm c}$ preserving a chosen chamber of
$W_{\rm r}$ is a complement.  On an even-sign packet it consists
of even products of the extra single-coordinate flips for the
nonempty size factors.  In particular
$W_\rho=W_{\rm r}\rtimes W_{\rm c}$.

\begin{lemma}[The presentation with strict transports]
\label{end:presentation}
Suppose the inducing simple module extends to its inertia group.
Then $\End_G(R_L^G\rho)$, or its opposite according to the action
convention, has the Hecke presentation of $W_{\rm r}$ with the
quadratic relations~\eqref{end:quadratic}, together with the
ordinary complement $W_{\rm c}$ acting on those generators by
conjugation.  It has a basis indexed by $W_\rho$.  The scalars
$\alpha_a$ agree for simple reflections joined by an odd Coxeter
bond and for generators conjugate under $W_{\rm c}$.
\end{lemma}

\begin{proof}
The extension supplies coefficient transports satisfying the group
law, with the prescribed Levi action on inner representatives.
For a gallery which crosses each reflecting hyperplane of
$W_{\rm r}$ at most once, use the gallery word theorem in the full
signed arrangement.  Braid moves preserve the change operator by
the minimal-gallery identity in Lemma~\ref{end:averaging}.  A
deletion cannot concern a repeated reflecting hyperplane, so it
is a pair of inverse changes at a noninertial wall.  Thus the
completed operator for a reduced gallery in the reflection
arrangement is its direct completed average, with no scalar
ambiguity.  This proves the braid and length-additive relations.
To compute a simple reflection at the base chamber, move through
noninertial walls to a chamber adjacent to its reflected chamber.
The changes used are invertible; conjugating the rank-one
calculation back gives~\eqref{end:quadratic}.

The direct averages also give the conjugation action of an
element preserving the chamber.  The operators so
obtained have single, distinct nonzero Mackey supports indexed by
$W_\rho$.  They are linearly independent and, by
\eqref{end:mackey}, span the endomorphism algebra.  Finally the
braid relations and invertibility identify the quadratics on
conjugate simple generators, proving the parameter assertions.
\end{proof}

\begin{lemma}[Extension of the inducing module]\label{end:extension}
Every module $\rho$ in~\eqref{end:rho} extends to $N_G(\mathbf L,\rho)$.
\end{lemma}

\begin{proof}
We construct the extension in three stages.  At the fine Levi,
an ordinary multiplicity-one constituent removes the scalar
cocycle.  We then obtain an invariant ordinary constituent for
the coarse Levi, whose multiplicity-one reduction supplies the
required modular extension.

Refine $\mathbf L$ to a Levi $\mathbf L'$ in which every chunk
has size one.  Let $\psi'$ be the ordinary cuspidal row at $s$
there, using $\chi_0$ and the size-one orbit rows.  Its inertia
quotient is a product of type $B$ groups and, on even-sign
packets, one type $D$ group; it is a reflection group without the
extra complement between sizes.  Iterating the positive split
single-cycle rule of Proposition~\ref{prop:single-cycle} identifies
its ordinary induction matrix with that of the ordinary
unipotent triangle series in $\mathbf C$.  The constituent
corresponding to the trivial relative Weyl character has
multiplicity one.  This is an ordinary character assertion, prior
to any modular endomorphism calculation.

Work first over an algebraically closed characteristic-zero
splitting field.  The coefficient intertwiners of $\psi'$ define
a scalar central extension of its inertia quotient: pairs
consisting of a normalizer element and an intertwiner are divided
by the natural pairs supplied by $L'$.  Normalize the lift of
each simple reflection to be an involution in this extension.
For two such lifts $\widetilde a,\widetilde b$ with Coxeter bond
$h$, the scalar
\[
 c_{ab}=(\widetilde a\widetilde b)^h
\]
is conjugate to its inverse by $\widetilde a$.  Hence
$c_{ab}=c_{ab}^{-1}$ and $c_{ab}\in\{1,-1\}$.
It is also, up to inversion, the multiplier between the two
completed braid words.  Indeed their underlying averages agree
by Lemma~\ref{end:averaging}; moving the coefficient transports
to the end leaves precisely the two words in these lifts.

On the multiplicity-one constituent of $R_{L'}^G\psi'$, every
completed operator acts by a nonzero scalar.  For an even bond
the two braid words have the same number of each generator, so
their scalar values coincide and $c_{ab}=1$.  For an odd bond,
changing the sign of either generator changes the multiplier.
The graph consisting of odd bonds in a classical Coxeter diagram
is a forest.  Choose signs successively along each tree to make
all its multipliers $1$.  This does not affect even bonds.
The Coxeter presentation now gives a section of the scalar
extension and therefore an ordinary extension of $\psi'$ to
its full inertia group.  Empty factors and $D_1$ require no
generators; $D_2$ has just the even commutation bond.

We transfer this extension to the coarse Levi.  Choose an $F$-stable
parabolic $\mathbf P'_L$ of $\mathbf L$ with Levi $\mathbf L'$
by ordering the size-one subchunks inside each chunk, and put
$P'_L=(\mathbf P'_L)^F$.  Every element of $W_\rho$ can be
lifted to preserve $\mathbf L'$, $\psi'$, and $\mathbf P'_L$: match subchunks in
order for a positive transport and in reverse order with all
signs changed for an orientation reversal.  For $j<h$ the latter
operation sends the positive linear root $e_j-e_h$ to
$e_{b+1-h}-e_{b+1-j}$, still positive.  The total type $D$ parity
is preserved, or the same residual sign compensation is used.
These are rational Weyl operations in $\mathbf C$ and give the
primal normalizer operations through the corresponding split
cocharacter positions.

Let $A$ be the subgroup of $N_G(\mathbf L,\rho)$ consisting of
elements that normalize $\mathbf L'$ and $\mathbf P'_L$ and
stabilize $\psi'$.  The preceding lifts show that $A$ maps onto
$W_\rho$, and the algebraic normalizer identities give
\[
 A\cap L\leq
 \bigl(N_{\mathbf L}(\mathbf P'_L)
       \cap N_{\mathbf L}(\mathbf L')\bigr)^F=L'.
\]
The strict action of $A$ on $\psi'$ and
conjugation on the group-algebra factor extend $R_{L'}^L\psi'$
to $LA=N_G(\mathbf L,\rho)$.  On $A\cap L$ this action is the natural
inner action, so it agrees with the $L$-module structure.
The ordinary constituent $\theta$ obtained by choosing the
regular extreme of Lemma~\ref{end:linear-cuspidals} in every
chunk and $\chi_0$ on the residual factor occurs once, by
ordinary general linear branching.  It is invariant: matched
chunks carry the same extreme and the residual cusp label is
fixed.  Its entire isotypic summand consequently inherits the
extension.

\smallskip\noindent\emph{Passage to the modular module.}
The extension of the ordinary constituent is now available.
Its reduction contains $\rho$ once, by
Lemma~\ref{end:linear-cuspidals} and the tensor product
description.  Reduce a stable lattice of the extended ordinary
module and choose a simple composition factor whose restriction
contains $\rho$.  Restriction of a simple module to a finite
normal subgroup is semisimple: the translates of any simple
submodule form a semisimple invariant sum, which is the entire
module.  Modular Clifford theory therefore makes this
restriction a sum of conjugates of $\rho$ with a common
multiplicity.  Here $\rho$ is invariant, and its total
multiplicity in the original reduction is one.  The chosen
factor restricts to $\rho$ itself, and is the required extension.
\end{proof}

\subsection{Rank-one parameters and the number of simple modules}

The extension lemma makes the presentation an untwisted one.
We now calculate its parameters and count its simple types.
The count must retain the single parity condition across all
chunk sizes in the even-sign case; imposing parity separately
at each size would give the wrong endpoint count.

For $v\in k^\times$, let $\mathcal H_v(\mathfrak S_n)$ denote
the type $A$ Hecke algebra with $(T-v)(T+1)=0$.  Its quantum
characteristic is the least positive $e$ for which
$1+v+\cdots+v^{e-1}=0$, or infinity if there is no such $e$.
Over $k$ it is $\ord(v)$ when $v\ne1$ and $p$ when $v=1$.
Its simple modules are indexed by $e$-regular partitions, namely
partitions in which no part occurs $e$ or more times
\cite{DipperJames}; see also \cite[Section~5 and Theorem~5.1]{DuRui}.

\begin{lemma}[The necessary parameters]\label{end:parameters}
After coherent rescaling of generators, the long parameter on
chunks of size $b$ in packet $i$ is $v_{i,b}=Q_i^b\in k^\times$.
On a type $B$ factor the two roots $z_1,z_2$ of the short-generator
quadratic satisfy
\begin{equation}\label{end:separation}
 z_1/z_2\notin \langle v_{i,b}\rangle.
\end{equation}
For $b>1$ the unrescaled short generator satisfies $T_0^2=1$.
On an even-sign packet, adjoining the single-coordinate flips
gives type $B$ factors with short generators of square $1$ and
the same long parameters.
\end{lemma}

\begin{proof}
Compute in the rank-one enlargement of a reflecting wall.
For a transposition of two equal-size chunks, choose the extra
$p$-parts in an ordinary lift $\psi^u$ to match on these chunks.
For a negative transposition orient the two chunks consistently
first.  The merged actual general linear block then has
centralizer $\GL_2(Q_i^b)$ on the matching full orbit.  Its two
ordinary induced constituents are the two Green partition rows,
each with multiplicity one and with degree ratio $Q_i^b$, up to
inversion; see \cite[Equation~(2.3.7)]{BrundanDipperKleshchev}.
The remaining factors are unchanged.

For a short wall with $b=1$, take trivial added $p$-part on that
chunk.  The ordinary lift is invariant there.  The split
single-cycle rule, applied at $su$ with all other added parts on
inactive chunks, identifies the rank-one matrix with the first
ordinary unipotent sign step above the residual triangle.
By Proposition~\ref{prop:label-degrees}, its two degrees have
ratio $Q_i^h$ for an integer $h$ in the orthogonal and
symplectic cases, or $q^{2t+1}$ at a unitary staircase
$t(t+1)/2$, again up to inversion.  Sign compensation on a
residual type $D$ triangle fixes the same inducing row.

The ordinary rank-one coefficient representation extends across
the reflection after scalar normalization.  Its completed
operator has zero trace, since its support is the nonidentity
double coset: in the induced-coset decomposition every diagonal
block is zero.  On the two multiplicity-one constituents, let
its roots be $z_1,z_2$ and their degrees be $D_1,D_2$.  Then
$D_1z_1+D_2z_2=0$, so the ratio of the roots is the negative
degree ratio, up to inversion.  Relation~\eqref{end:quadratic}
also gives $z_1z_2=-1$.

Choose a lattice stable under the finite rank-one inertia group,
enlarging the coefficient field to contain the square root of
$r$ used in~\eqref{end:average}.  The averaging denominators and
normalization scalars are units, so the completed operator preserves
the induced lattice.  Its eigenvalues are integral and have product
$-1$, hence are units.  Their sum $\alpha_a$ is integral, so
$T_a^{-1}=T_a-\alpha_a$ also preserves the lattice.
The restriction of the coefficient lattice reduces
irreducibly to $\rho$, and its involutive transport reduces to
our chosen one up to sign.  Thus the same root ratios give the
modular quadratic, even if its two long roots coincide.

For a short wall with $b>1$, instead choose the ordinary lift
noninvariant at that wall.  In the orthogonal and symplectic
cases its regular $p$-element cannot be conjugate to its inverse:
already modulo $p$, $-1\notin\langle Q_i\rangle$, since the latter
group has odd order.  In the unitary case its paired transform
would require $-q\in\langle q^2\rangle$, which is excluded by
$\ord_p(-q)=2d_i$ and $\ord_p(q^2)=d_i$.
To check the marking, choose the dual wall representative
$n\in\mathbf C^F$, so $n(s)=s$.  Any conjugacy of $su$ and
$s n(u)$ in $\mathbf L^{*F}$ must match their $p'$-parts, forcing
the conjugator into $\mathbf C_L^F$.  Its active factor is
$\GL_b(Q_i)$; in a degree-two packet this is $\GL_b(q^2)$,
so eigenvalue conjugacy uses $q^2$-powers.  The preceding
exclusion therefore applies in that case as well.
The ordinary reverse composition of uncompleted
averages is therefore the identity by rank-one Mackey.  Reducing
it and using the involutive modular transport gives $T_0^2=1$.
Different walls may use different ordinary lifts; no simultaneous
extension of these lifts is required.

For $b=1$ the short-root ratio is $-Q_i^h$ in the orthogonal and
symplectic cases.  It is outside $\langle Q_i\rangle$ by odd
order.  In the unitary case the ratio is
$-q^{2t+1}=(-q)^{2t+1}$, outside the subgroup of even powers
$\langle q^2\rangle$.  For $b>1$, $d_i\mid b$, so
$v_{i,b}=1$ in $k$; the two short roots are opposite and distinct
because $p$ is odd.  This proves~\eqref{end:separation}.

The rescalings giving $(T-v_{i,b})(T+1)=0$ can be chosen on
conjugacy classes of simple generators by
Lemma~\ref{end:presentation}.  This includes the two ends of
$D_2$ when a complement interchanges them; otherwise they may
be rescaled independently.  Adjoining a single-coordinate flip
to a type $D$ factor gives its standard type $B$ presentation
with short generator of square $1$.  In the empty-residual case
this also follows directly from the coordinate signed
permutation presentation.  The convention is valid for $D_1=1$.
\end{proof}

We use the following specific separated-parameter theorem.  If
the cyclotomic Hecke algebra with type $A$ parameter $v$ has
two short parameters $z_1,z_2$ and
$v^a z_1-z_2$ is a unit for every integer $a$ with $-n<a<n$,
then it is Morita equivalent to
\begin{equation}\label{end:dm}
 \bigoplus_{a=0}^n
  \mathcal H_v(\mathfrak S_a)\otimes
  \mathcal H_v(\mathfrak S_{n-a}).
\end{equation}
This is the separated-parameter theorem of Du--Rui
\cite[Theorem~4.14]{DuRui}, also the two-singleton case of
\cite[Theorem~1.1 and Corollary~1.4]{DipperMathas}; its base
ring can be a field of positive characteristic.  Our parameters
are nonzero, and~\eqref{end:separation} verifies its hypothesis
for every integer $a$, uniformly in $n$.

\begin{lemma}[Counting the even subalgebra]\label{end:even-algebra}
For one even-sign packet with at least one chunk, the simple
modules of its endomorphism algebra are indexed by the following
rule.  For every size $b$ choose an ordered pair of
$e_b$-regular partitions of total size $r_{i,b}$, where
\begin{equation}\label{end:quantum-e}
 e_b=\begin{cases}
 d_i,&b=1\text{ and }d_i>1,\\
 p,&\text{otherwise}.
 \end{cases}
\end{equation}
Identify these choices under simultaneous exchange of the two
components for all sizes.  Each nonfixed orbit contributes one
simple, and each fixed choice contributes two.  If there are no
chunks, there is one simple.
\end{lemma}

\begin{proof}
Let $\mathcal B$ be the tensor product, over the nonempty sizes,
of the type $B$ Hecke algebras with short generator of square
$1$ from Lemma~\ref{end:parameters}.  Give every short generator
odd degree and every long generator even degree.  The algebra
of our packet is the total-even subalgebra $\mathcal B_0$.
Indeed its reflection subgroup is the product of the type $D$
groups; its complement consists of even products of the extra
flips, precisely as in~\eqref{end:total-parity}.

The separated equivalence~\eqref{end:dm} and the type $A$ simple
module theorem label the simples of $\mathcal B$ by the ordered
pairs in the statement.  Formula~\eqref{end:quantum-e} follows
because $Q_i$ has order $d_i$, while $Q_i^b=1$ for all the other
allowed sizes.  In particular parameter $1$ has quantum
characteristic $p$, not $1$.

The grading automorphism $\tau$ of $\mathcal B$ negates all
short generators and fixes all long generators.  On the
separated labels it exchanges the two components at every
size.  To check the action on their type $A$ labels as well as
on their multiplicities, use the affine presentation
\[
 X_1=T_0,\qquad X_{j+1}=v^{-1}T_jX_jT_j.
\]
The $X_j$ commute, and $\tau$ negates them while leaving $T_j$
unchanged.  Their generalized weights lie in the two disjoint
strings $z_1v^{\Z}$ and $z_2v^{\Z}$, here with $z_2=-z_1$.
The corner with all strands of each string grouped together
has the two type $A$ factors in~\eqref{end:dm}.  Negation
exchanges its two groups.

There is no additional involution on either type $A$ factor.
For completeness, the Bernstein intertwiners
\[
 \Phi_j=T_j-\frac{v-1}{1-X_j/X_{j+1}}
\]
interchange the $X_j$ weights and satisfy the braid relations
where the indicated differences are invertible.  These
relations follow by substitution in the affine Hecke
relations.  Between different strings they are invertible,
since both equal-weight and $v$-multiple-weight coincidences
are excluded.  The product which swaps the two grouped blocks
preserves the order inside each block.  The intertwiner braid
relations show that it conjugates a within-block $\Phi_j$ to
the corresponding within-block $\Phi_h$, and it carries the
displayed rational correction to the corresponding correction.
It therefore conjugates $T_j$ to $T_h$.  This calculation may
first be made in the Laurent localization.  After the
within-block corrections cancel, only differences between the
two strings have been inverted; the affine basis theorem
therefore gives the same identity in the separated corner.
Thus the two type $A$ modules are simply exchanged, with their
partition labels unchanged.

The action on the separated partition labels is now determined.
It remains to restrict from the full signed algebra to its
total-even subalgebra.
We have $\mathcal B=\mathcal B_0\oplus\mathcal B_0u$ for an
invertible odd short generator $u$ with $u^2=1$.  Restriction
of a simple $\mathcal B$-module to $\mathcal B_0$ is semisimple:
the sum of a simple submodule and its $u$-translate is a
nonzero $\mathcal B$-submodule.  The usual index-two Clifford
argument now applies, since $2$ is invertible in $k$.  A
simple and its distinct $\tau$-twist restrict to the same
simple module.  A simple fixed by $\tau$ has a normalized
twisting intertwiner of square $1$; its two eigenspaces give
two distinct simple restrictions.  These exhaust the simple
$\mathcal B_0$-modules by induction from $\mathcal B_0$.
This gives exactly the stated rule.  With no chunks there is
no odd unit and the algebra is just $k$, explaining the
separate convention.
\end{proof}

\begin{lemma}[Exact total count]\label{end:count}
Sum the numbers of simple endomorphism-algebra types over all
size choices~\eqref{end:size-sum}.  The result is $|\mathcal T|$.
\end{lemma}

\begin{proof}
For an integer $e\geq2$ put
\[
 R_e(x)=\prod_{j\geq1}(1+x^j+\cdots+x^{(e-1)j}),
 \qquad
 P(x)=\prod_{j\geq1}(1-x^j)^{-1}.
\]
The first series counts $e$-regular partitions, the second all
partitions.  If $d>1$, expansion of each part multiplicity
first modulo $d$ and then in base $p$ gives
\begin{equation}\label{end:digits}
 R_d(x)\prod_{a\geq0}R_p(x^{dp^a})=P(x).
\end{equation}
For $d=1$ the corresponding identity is
$\prod_{a\geq0}R_p(x^{p^a})=P(x)$.  These are identities of
formal series: at each degree only finitely many factors
matter, and the products telescope using
$R_e(x)=P(x)/P(x^e)$.  Equivalently, a partition $\lambda_b$
chosen at size $b$ contributes $b$ copies of each of its
parts to the combined partition.  Uniqueness of the digits
is precisely uniqueness of this decomposition.

For an ordinary signed packet, the two partition components
are independent.  Squaring~\eqref{end:digits} thus counts all
ordered bipartitions of $m_i$, exactly the ordinary type $B$
triangle labels.  The construction commutes with exchanging
the two components.  Hence on an even-sign packet the rule
of Lemma~\ref{end:even-algebra} counts unordered bipartitions,
with two labels on equal pairs, exactly the type $D$ rule.
For $m>0$, if $B(m)$ is the number of ordered bipartitions and
$F(m)$ the number of equal pairs, the count is
\[
 \frac{B(m)-F(m)}2+2F(m)=\frac{B(m)+3F(m)}2,
 \qquad
 F(m)=\begin{cases}P_{m/2},&m\text{ even},\\0,&m\text{ odd},\end{cases}
\]
where $P_a$ is the number of partitions of $a$.  At $m=0$
both sides of the representation-theoretic count use one
label, instead of applying this positive-rank formula.
In particular $D_1$ gives one label and $D_2$ gives four.
The latter also follows directly from its two rank-one
Hecke factors: their parameter has odd order, or is $1$ in
odd characteristic, and is therefore not $-1$.

If several sizes each have just one chunk, their reflection
groups $D_1$ are trivial but their total-even complement is
$C_2^{t-1}$ for $t$ sizes.  The simultaneous exchange rule
gives $2^{t-1}$ labels as required.  Thus the complement
has not been lost in any small-rank case.  Finally the
packet counts multiply, proving the lemma.
\end{proof}

\subsection{Exhaustion and the projective norm estimate}

The total endomorphism-algebra count is now $|\mathcal T|$.
By Lemma~\ref{end:head}, this is the number of distinct simple
modules in our disjoint heads.  We finish the exhaustion argument
by placing all these heads in the endpoint block union.  We then
bound their projective covers by inducing projectives from the
bounded chunks and the central-defect residual block.

\begin{corollary}[Exhaustion of the endpoint simples]
\label{end:exhaustion}
Every simple module in the block union of $\mathcal T$
is a quotient of $R_L^G\rho$ for one of the pairs
constructed in~\eqref{end:rho}.
\end{corollary}

\begin{proof}
The Brauer class of $\rho$ is an integral combination of
partition rows at the marked parameter of $L$, with
residual row $\chi_0$.  Every one of these ordinary rows
occurs in $R_{L'}^L\psi'$, by ordinary linear branching.
Transitivity and Proposition~\ref{prop:single-cycle}
therefore show that their inductions use only rows of
$\mathcal T$.  Thus $R_L^G\rho$ has all its composition
factors in the block union in question; the integral
combination is enough for this support assertion.

Lemmas~\ref{end:presentation}--\ref{end:count} count the
simple types of all the relevant endomorphism algebras.
Lemma~\ref{end:head} identifies this with the count for their heads.
These heads are disjoint for different size choices, by
Lemmas~\ref{end:inertia} and~\ref{end:head}.  The resulting total is
$|\mathcal T|$.  Because $\mathcal T$ is an integral
basic set of the entire block union, this is exactly its
number of simple modules.  The supported heads therefore
exhaust all of them.
\end{proof}

\begin{proof}[Proof of Proposition~\ref{prop:triangle-cartan}]
Let $S$ be one of the endpoint simples.  By
Corollary~\ref{end:exhaustion} it is a quotient of
$R_L^G\rho$.  Let $P_\rho$ be the projective cover of
$\rho$.  Exactness and preservation of projectives give
a projective module $R_L^G P_\rho$ mapping onto $S$.
After projection to the endpoint block union, it has the
projective cover $P_S$ as a direct summand.

We bound the ordinary coordinates of $P_\rho$ before
inducing.  There are at most $\sum_i m_i$ new chunks,
each of size at most this sum, so each actual new linear
factor has rank at most $2\sum_i m_i$.  Their Sylow
$p$-orders are bounded as well.  Explicitly, for
$e_q=\ord_p(q)$ and $c_q=v_p(q^{e_q}-1)$, lifting the
exponent gives
\begin{equation}\label{end:gl-order}
 v_p(|\GL_N(q)|)
 =c_q\left\lfloor\frac N{e_q}\right\rfloor
  +v_p\!\left(\left\lfloor\frac N{e_q}\right\rfloor!\right).
\end{equation}
For $Q_i=q$ or $q^2$, $c_q$ is controlled by
$v_p(Q_i^{d_i}-1)$; the possible factor $2$ does not
change the valuation since $p$ is odd.  For unitary
chunks apply the same formula over $Q_i=q^2$.
The geometric bound of Theorem~\ref{thm:geometric-cartan}
therefore bounds their Cartan entries, since both rank
and defect order are bounded.  Lemma~\ref{end:residual}
bounds the residual block; frozen factors have defect
zero.  The actual product decomposition and central
transfer consequently give a bound $C_0$ for the
diagonal Cartan entry of $P_\rho$.  Thus its ordinary
projective-character coefficients are at most
$\sqrt{C_0}$, and its squared ordinary norm is at most
$C_0$.  The number $N_0$ of its nonzero ordinary
coordinates is bounded, either by this integral norm
bound or by the bounded-defect row bound.

Only ordinary coordinates at the marked $p'$-parameter
of $L$ contribute to the $s$-projection of its induction.
Indeed an ordinary coordinate with nontrivial extra
$p$-part retains that nontrivial part under induction
and cannot belong to the series at $s$.
The selected Levi block union fixes the marked
$p'$-parameter class, so other Levi classes fusing to
$s$ are not included.  The residual coordinate is
$\chi_0$ by Lemma~\ref{end:residual}.

Every remaining row $\theta$ of $L$ occurs with positive
integer multiplicity in $R_{L'}^L\psi'$.  Positivity of
ordinary split induction and transitivity give the
coefficientwise inequality
\[
 R_L^G\theta\ \leq\ R_{L'}^G\psi'.
\]
The ordinary positive transfer of
Proposition~\ref{prop:single-cycle} identifies the
latter multiplicities with ordinary relative Weyl
multiplicities in the triangle series.  They are
bounded, for example, by
\[
 W_0=\prod_i 2^{m_i}m_i!.
\]
This estimate concerns coordinate multiplicities; it
does not bound ordinary character degrees.
It follows that every $\mathcal T$-coordinate of the
projective character of $R_L^G P_\rho$ is at most
$N_0\sqrt{C_0}W_0$.  Since $P_S$ is its projective
summand and all ordinary projective coefficients are
nonnegative, the same coordinate bound holds for
$P_S$.

Both the number of basic coordinates in $\mathcal T$
and the number of ordinary irreducible characters in
its block union are bounded, as are the block defect orders, by
Proposition~\ref{prop:runner-reduction}.  Apply
Lemma~\ref{lem:row-projection} to recover a uniform
bound for the full ordinary norm of $P_S$ from this
projected bound.  Its squared norm is its diagonal
Cartan entry, and Cauchy--Schwarz bounds every off-diagonal
entry by the corresponding diagonal bounds.  This
proves the proposition, including the all-zero-rank
case, in which the residual central-defect argument
alone applies.
\end{proof}
